\subsection{约化范数与约化迹的性质}

命题 4. —— 设 L 是 K 的一个扩张，a 是中心单代数 A 的一个元素。我们有关系式

\[
\operatorname{Pcrd} _ {\mathrm{A} _ {(\mathrm{L})} / \mathrm{L}} (1 \otimes a; \mathrm{X}) = \operatorname{Pcrd} _ {\mathrm{A} / \mathrm{K}} (a; \mathrm{X}),\tag{28}
\]

\[
\mathrm{Trd} _ {\mathrm{A} _ {(\mathrm{L})} / \mathrm{L}} (1 \otimes a) = \mathrm{Trd} _ {\mathrm{A} / \mathrm{K}} (a),\tag{29}
\]

\[
\mathrm{Nrd} _ {\mathrm{A} _ {(\mathrm{L})} / \mathrm{L}} (1 \otimes a) = \mathrm{Nrd} _ {\mathrm{A} / \mathrm{K}} (a).\tag{30}
\]

（``在标量扩张下的不变性''）。

等式 (28) 的两边有相同的 \(n\) 次幂，这是由定义（VIII, 第 340 页的公式 (20)）以及关系式

\[
\mathrm{Pc} _ {\mathrm{A} _ {(\mathrm{L})} / \mathrm{L}} (1 \otimes a; \mathrm{X}) = \mathrm{Pc} _ {\mathrm{A} / \mathrm{K}} (a; \mathrm{X})
\]

（III, §9, No.~3, 第 544 页，公式 (21)）得到的。因此等式 (28) 由 VIII, 第 339 页的引理 2 得出。公式 (29) 与 (30) 由 (28)、(21) 与 (22) 得出。

推论 1. —— 设 L 是 K 的一个扩张。设 V 是 L 上一个维数为 n 的向量空间，\(\theta\) 是一个 K-代数同态（从 A 到 \(\operatorname{End}_{\mathrm{L}}(\mathrm{V})\) 的）。对 A 的每个元素 a，我们有

\[
\mathrm{Pcrd} _ {\mathrm{A/K}} (a; \mathrm{X}) = \chi_ {\theta (a)} (\mathrm{X}),\tag{31}
\]

\[
\mathrm{Trd} _ {\mathrm{A/K}} (a) = \mathrm{Tr} (\theta (a)),\tag{32}
\]

\[
\mathrm{Nrd} _ {\mathrm{A/K}} (a) = \det (\theta (a)).\tag{33}
\]

设 \(\theta'\) 是从 \(\mathrm{A}_{(\mathrm{L})}\) 到 \(\mathrm{End}_{\mathrm{L}}(\mathrm{V})\) 的 L-代数同态，它满足 \(\theta'(\lambda \otimes a) = \lambda \theta(a)\)（对 \(\lambda \in \mathrm{L}\) 与 \(a \in \mathrm{A}\)）。由 VIII, 第 222 页的推论 2，代数 \(\mathrm{A}_{(\mathrm{L})}\) 是单的；因此同态 \(\theta'\) 是单射。但是体 L 上的代数 \(\mathrm{A}_{(\mathrm{L})}\) 与 \(\mathrm{End}_{\mathrm{L}}(\mathrm{V})\) 有相同的次数，都等于 \(n^2\)，所以 \(\theta'\) 是一个同构。于是推论 1 由命题 4 与上面的例 2 得出。

推论 2. —— 设 \(a\) 与 \(a'\) 是 A 的元素，\(\lambda\) 是 K 的一个元素。我们有关系式

\[
\mathrm{Pcrd} _ {\mathrm{A/K}} (a; a) = 0,\tag{34}
\]

\[
\mathrm{Pcrd} _ {\mathrm{A/K}} (\lambda a; \lambda \mathrm{X}) = \lambda^ {n} \mathrm{Pcrd} _ {\mathrm{A/K}} (a; \mathrm{X}),\tag{35}
\]

\[
\mathrm{Trd} _ {\mathrm{A/K}} (a + a ^ {\prime}) = \mathrm{Trd} _ {\mathrm{A/K}} (a) + \mathrm{Trd} _ {\mathrm{A/K}} (a ^ {\prime}),\tag{36}
\]

\[
\mathrm{Trd} _ {\mathrm{A/K}} (\lambda a) = \lambda \mathrm{Trd} _ {\mathrm{A/K}} (a),
\]

\[
\mathrm{Trd} _ {\mathrm{A/K}} (a a ^ {\prime}) = \mathrm{Trd} _ {\mathrm{A/K}} (a ^ {\prime} a),\tag{37}
\]

\[
\mathrm{Nrd} _ {\mathrm{A/K}} (a a ^ {\prime}) = \mathrm{Nrd} _ {\mathrm{A/K}} (a) \cdot \mathrm{Nrd} _ {\mathrm{A/K}} (a ^ {\prime}),\tag{38}
\]

\[
\mathrm{Nrd} _ {\mathrm{A/K}} (\lambda a) = \lambda^ {n} \mathrm{Nrd} _ {\mathrm{A/K}} (a),
\]

\[
\operatorname{Trd} _ {\mathrm{A} / \mathrm{K}} (1) = n, \operatorname{Nrd} _ {\mathrm{A} / \mathrm{K}} (1) = 1.\tag{39}
\]

由于 A 是中心单的并且在 K 上的约化次数为 n，存在 K 的一个扩张 L 以及 L 上一个维数为 n 的向量空间 V，使得 \(\mathrm{A}_{(\mathrm{L})}\) 同构于代数 \(\mathrm{End}_{\mathrm{L}}(\mathrm{V})\)（VIII, 第 252 页，定理 1）。于是推论 2 由推论 1 以及自同态的迹与行列式的性质得出。特别地，公式 (34) 由凯莱--哈密顿定理得出（III, §8, No.~11, 第 541 页，命题 20，另参看 VIII, 第 338 页，注 2）。

推论 3. —— 设 \(A^{o}\) 是 A 的反代数。对 A 中的每个 a，我们有

\[
\mathrm{Pcrd} _ {\mathrm {A^ {\circ} /K}} (a; \mathrm{X}) = \mathrm{Pcrd} _ {\mathrm{A/K}} (a; \mathrm{X}),\tag{40}
\]

\[
\mathrm{Trd} _ {\mathrm {A^{\circ} /K}} (a) = \mathrm{Trd} _ {\mathrm{A/K}} (a),\tag{41}
\]

\[
\mathrm{Nrd} _ {\mathrm {A^{\circ} /K}} (a) = \mathrm{Nrd} _ {\mathrm{A/K}} (a).\tag{42}
\]

选取 K 的一个扩张 L、L 上一个维数为 n 的向量空间 V，以及一个同态 \(\theta\)（从 A 到 \(\operatorname{End}_{\mathrm{L}}(\mathrm{V})\) 的）（VIII, 第 252 页，定理 1）。设 \(V^{*}\) 是 V 的对偶向量空间。把 A 的元素 a 映到自同态 \({}^{t}\theta(a)\)（\(V^{*}\) 的）的映射是一个 K-代数同态（从 \(A^{o}\) 到 \(\operatorname{End}_{\mathrm{L}}(\mathrm{V}^{*})\) 的）。于是推论 3 由推论 1 与 III, §8, No.~4, 第 528 页的推论 3 得出。

因此，a 的迹、范数与特征多项式，不论我们把 a 看作 A 的元素还是 \(A^{\circ}\) 的元素，都是一样的。当 A 不假定为中心单时，这一性质并不总成立（III, §9, 第 644 页，习题 1）。

\pandocbounded{\includegraphics[keepaspectratio]{images/part002_d5c76b31.jpg}}

命题 5. —— 对 A 中的任何 x，令 \(t_{x}\) 是 A 上的线性形式 \(y \mapsto \operatorname{Trd}_{\mathrm{A/K}}(xy)\)。

\begin{enumerate}
\def\labelenumi{\alph{enumi})}
\item
  映射 \(t: x \mapsto t_x\) 是一个 (A, A)-双模同构（从 A 到它的对偶 Hom\(_{K}\) (A, K) 的）。
\item
  设 \(h\) 是 \(A\) 上的一个线性形式。下列性质等价：
\end{enumerate}

\begin{enumerate}
\def\labelenumi{(\roman{enumi})}
\item
  存在一个元素 \(\lambda\)（\(K\) 的），使得 \(h(x) = \lambda \operatorname{Trd}_{A / K}(x)\)（对每个 \(x \in A\)）。
\item
  我们有 \(h(xy) = h(yx)\)（对 A 中所有的 x, y）。
\end{enumerate}

回顾（II, §1, No.~14, 第 225--226 页），\(A^{*} = \operatorname{Hom}_{K}(A, K)\) 上的 (A, A)-双模结构由关系式

\[
\langle a t b, c \rangle = \langle t, b c a \rangle\tag{43}
\]

（对 \(a, b, c \in A\) 与 \(t \in A^{*}\)）定义。特别地，对 A 中的每个 x，我们有

\[
\begin{array}{c} \langle a t _ {x} b, c \rangle = \langle t _ {x}, b c a \rangle = \mathrm{Trd} _ {\mathrm{A/K}} (x b c a), \\ \langle t _ {a x b}, c \rangle = \mathrm{Trd} _ {\mathrm{A/K}} (a x b c), \end{array}
\]

并且这两个元素由 VIII, 第 342 页的公式 (37) 可知相等。因此我们有 \(at_x b = t_{axb}\)，这就意味着 \(t\) 是一个 (A, A)-双模同态（从 A 到 A* 的）。

我们选取体 K 的一个扩张 L 以及一个同构 \(\theta\)（从 L-代数 \(\mathrm{A}_{(\mathrm{L})}\) 到矩阵代数 \(\mathbf{M}_{n}(\mathrm{L})\) 的）（VIII, 第 252 页，定理 1）。我们

把向量空间 \((\mathrm{A}^{*})_{(\mathrm{L})}\) 与体 L 上的向量空间 \(\mathrm{A}_{(\mathrm{L})}\) 的对偶等同起来。按照这些约定，由 VIII, 第 341 页的命题 4，我们有

\[
\operatorname{Trd} _ {\mathrm{A} _ {(\mathrm{L})} / \mathrm{L}} = 1 _ {\mathrm{L}} \otimes \operatorname{Trd} _ {\mathrm{A} / \mathrm{K}}.\tag{44}
\]

设 \(t_{(L)}\) 是从 \(A_{(L)}\) 到 \(A_{(L)}^{*}\) 的 L-线性映射（由 t 经标量扩张得到的）；由公式 (44) 与推论 1，我们有

\[
\langle t _ {(\mathrm{L})} (x), y \rangle = \operatorname{Trd} _ {\mathrm{A} _ {(\mathrm{L})} / \mathrm{L}} (x y) = \operatorname{Tr} (\theta (x) \theta (y))\tag{45}
\]

（\(x, y\) 在 \(\mathrm{A}_{(\mathrm{L})}\) 中）。由 II, §10, No.~11, 第 358 页的命题 7，映射 \(t_{(\mathrm{L})}\) 是双射；由此推出 \(t\) 是双射。我们证明了 a)。

设 \(h\) 属于 \(\mathrm{A}^*\)；由上面的结果，存在一个元素 \(a\)（\(\mathrm{A}\) 的），使得 \(h\) 等于 \(t_a\)。由 a)，我们有

\[
h (x y) - h (y x) = t _ {a} (x y - y x) = t _ {a x} (y) - t _ {x a} (y).
\]

因此，关系式`` \(h(xy) = h(yx)\)（对 A 中的 \(x,y\)）''等价于`` \(t_{ax-xa} = 0\)（对每个 \(x \in \mathrm{A}\)）''，而由证明中的 a) 部分，这意味着 \(a\) 属于 A 的中心 K。于是结合公式 (36)，这就证明了 b)。

推论. —— A 中由形如 xy - yx 的元素（其中 x 与 y 取遍 A）所生成的线性子空间是一个超平面，即非零线性形式 \(Trd_{A/K}\) 的核。

注. —— 由 VIII, 第 340 页的公式 (23)，我们有

\[
\mathrm{Tr} _ {\mathrm{A/K}} (a) = n \mathrm{Trd} _ {\mathrm{A/K}} (a)
\]

（对每个 \(a \in A\)）。如果体 K 的特征等于 0 或一个不整除 n 的素数 p，那么在命题 5 中我们可以用迹来代替约化迹。另一方面，如果 K 的特征是一个整除 n 的素数，那么我们有 \(\operatorname{Tr}_{\mathrm{A}/\mathrm{K}}(a) = 0\)（对每个 \(a \in A\)）。

